% GENERATED FILE. Edit canonical lessons, then run npm run generate:books.
% canonical-source-hash: fnv1a64:1356cff51d76fbbd
% canonical-lessons: GU-C142-kahevu, GU-C142-jivvu, GU-C142-kharcvu, GU-C142-bacavvu, GU-C142-vavvu

\chapter{To say, To live, To spend, To save, To sow}
\label{ch:gu-to-say-to-live-to-spend-to-save-to-sow}
\hlchaptermodality{142}

\begin{chapteropening}
\textbf{By the end of this chapter:} \emph{I can say to say, to live, to spend, to save and to sow.}

Complete the last lesson of chapter 142: I can say to say, to live, to spend, to save and to sow.
\end{chapteropening}

\section[\texorpdfstring{\gu{કહેવું} (kahevũ)}{કહેવું (kahevũ)}]{\gu{કહેવું} (kahevũ) — to say, to tell}

\label{lesson:GU-C142-kahevu}



\begin{quote}
Before the new one: say the Gujarati for to drive, then the Gujarati for to count.
\end{quote}

\subsection*{You'll want to know: \gu{કહેવું}}
\textbf{\gu{કહેવું}} (\emph{kahevũ}) — \textquotedblleft{}to say, to tell\textquotedblright{}.

\begin{cousinweb}
\textbf{\gu{કહેવું}} (\emph{kahevũ}) is a tadbhava, a word inherited from Sanskrit and reshaped by sound change. It comes from \emph{kathayati}, \textquotedblleft{}converses with, describes\textquotedblright{} (CDIAL 2703).
\end{cousinweb}

\begin{grammarlens}[title={What you've built}]
One more word for this chapter.
\end{grammarlens}

\subsection*{Your turn}
\begin{itemize}
\raggedright
  \item \emph{Say it:} \emph{kahevũ}
  \item \emph{Say it:} \emph{kahevũ}, once more
  \item \emph{Say it:} say \emph{kahevũ}
  \item \emph{Recall:} say the Gujarati for to drive, then the Gujarati for to count, then say \emph{kahevũ} again
\end{itemize}

\begin{tcolorbox}[breakable,colback=teal!4,colframe=teal!35!black,title={Before you move on}]
What does \gu{કહેવું} mean? (\textquotedblleft{}To say, to tell\textquotedblright{}.) Say it once more.
\end{tcolorbox}
\section[\texorpdfstring{\gu{જીવવું} (jīvvũ)}{જીવવું (jīvvũ)}]{\gu{જીવવું} (jīvvũ) — to live}

\label{lesson:GU-C142-jivvu}



\begin{quote}
Before the new one: say the Gujarati for to count, then the Gujarati for to say.
\end{quote}

\subsection*{You'll want to know: \gu{જીવવું}}
\textbf{\gu{જીવવું}} (\emph{jīvvũ}) — \textquotedblleft{}to live\textquotedblright{}.

\begin{cousinweb}
\textbf{\gu{જીવવું}} (\emph{jīvvũ}) is a tadbhava, a word inherited from Sanskrit and reshaped by sound change. It comes from \emph{jīvati}, \textquotedblleft{}is alive\textquotedblright{} (CDIAL 5241).
\end{cousinweb}

\begin{grammarlens}[title={What you've built}]
One more word for this chapter.
\end{grammarlens}

\subsection*{Your turn}
\begin{itemize}
\raggedright
  \item \emph{Say it:} \emph{jīvvũ}
  \item \emph{Say it:} \emph{jīvvũ}, once more
  \item \emph{Say it:} say \emph{jīvvũ}
  \item \emph{Recall:} say the Gujarati for to count, then the Gujarati for to say, then say \emph{jīvvũ} again
\end{itemize}

\begin{tcolorbox}[breakable,colback=teal!4,colframe=teal!35!black,title={Before you move on}]
What does \gu{જીવવું} mean? (\textquotedblleft{}To live\textquotedblright{}.) Say it once more.
\end{tcolorbox}
\section[\texorpdfstring{\gu{ખર્ચવું} (kharcvũ)}{ખર્ચવું (kharcvũ)}]{\gu{ખર્ચવું} (kharcvũ) — to spend}

\label{lesson:GU-C142-kharcvu}



\begin{quote}
Before the new one: say the Gujarati for to say, then the Gujarati for to live.
\end{quote}

\subsection*{You'll want to know: \gu{ખર્ચવું}}
\textbf{\gu{ખર્ચવું}} (\emph{kharcvũ}) — \textquotedblleft{}to spend\textquotedblright{}.

\begin{cousinweb}
\textbf{\gu{ખર્ચવું}} (\emph{kharcvũ}) is built on \emph{kharc}, \textquotedblleft{}expense\textquotedblright{}, usually traced through Persian \emph{kharch} to Arabic \emph{kharj}, \textquotedblleft{}expenditure; tribute\textquotedblright{}. The Arabic root is \emph{kh-r-j} (Lane).
\end{cousinweb}

\begin{grammarlens}[title={What you've built}]
One more word for this chapter.
\end{grammarlens}

\subsection*{Your turn}
\begin{itemize}
\raggedright
  \item \emph{Say it:} \emph{kharcvũ}
  \item \emph{Say it:} \emph{kharcvũ}, once more
  \item \emph{Say it:} say \emph{kharcvũ}
  \item \emph{Recall:} say the Gujarati for to say, then the Gujarati for to live, then say \emph{kharcvũ} again
\end{itemize}

\begin{tcolorbox}[breakable,colback=teal!4,colframe=teal!35!black,title={Before you move on}]
What does \gu{ખર્ચવું} mean? (\textquotedblleft{}To spend\textquotedblright{}.) Say it once more.
\end{tcolorbox}
\section[\texorpdfstring{\gu{બચાવવું} (bacāvvũ)}{બચાવવું (bacāvvũ)}]{\gu{બચાવવું} (bacāvvũ) — to save}

\label{lesson:GU-C142-bacavvu}



\begin{quote}
Before the new one: say the Gujarati for to live, then the Gujarati for to spend.
\end{quote}

\subsection*{You'll want to know: \gu{બચાવવું}}
\textbf{\gu{બચાવવું}} (\emph{bacāvvũ}) — \textquotedblleft{}to save\textquotedblright{}.

\begin{cousinweb}
\textbf{\gu{બચાવવું}} (\emph{bacāvvũ}) goes back to Sanskrit \emph{vañcati}, \textquotedblleft{}moves; escapes from\textquotedblright{} (CDIAL 11208). CDIAL marks this form as taken from Hindi; the native Gujarati form is \emph{vāṁcāvvũ}.
\end{cousinweb}

\begin{grammarlens}[title={What you've built}]
One more word for this chapter.
\end{grammarlens}

\subsection*{Your turn}
\begin{itemize}
\raggedright
  \item \emph{Say it:} \emph{bacāvvũ}
  \item \emph{Say it:} \emph{bacāvvũ}, once more
  \item \emph{Say it:} say \emph{bacāvvũ}
  \item \emph{Recall:} say the Gujarati for to live, then the Gujarati for to spend, then say \emph{bacāvvũ} again
\end{itemize}

\begin{tcolorbox}[breakable,colback=teal!4,colframe=teal!35!black,title={Before you move on}]
What does \gu{બચાવવું} mean? (\textquotedblleft{}To save\textquotedblright{}.) Say it once more.
\end{tcolorbox}
\section[\texorpdfstring{\gu{વાવવું} (vāvvũ)}{વાવવું (vāvvũ)}]{\gu{વાવવું} (vāvvũ) — to sow}

\label{lesson:GU-C142-vavvu}



\begin{quote}
Before the new one: say the Gujarati for to spend, then the Gujarati for to save.
\end{quote}

\subsection*{You'll want to know: \gu{વાવવું}}
\textbf{\gu{વાવવું}} (\emph{vāvvũ}) — \textquotedblleft{}to sow\textquotedblright{}.

\begin{cousinweb}
\textbf{\gu{વાવવું}} (\emph{vāvvũ}) is a tadbhava, a word inherited from Sanskrit and reshaped by sound change. It comes from \emph{vāpayati}, \textquotedblleft{}sows\textquotedblright{} (CDIAL 11525).
\end{cousinweb}

\begin{grammarlens}[title={What you've built}]
One more word for this chapter.
\end{grammarlens}

\subsection*{Your turn}
\begin{itemize}
\raggedright
  \item \emph{Say it:} \emph{vāvvũ}
  \item \emph{Say it:} \emph{vāvvũ}, once more
  \item \emph{Say it:} say \emph{vāvvũ}
  \item \emph{Recall:} say the Gujarati for to spend, then the Gujarati for to save, then say \emph{vāvvũ} again
\end{itemize}

\begin{tcolorbox}[breakable,colback=teal!4,colframe=teal!35!black,title={Before you move on}]
What does \gu{વાવવું} mean? (\textquotedblleft{}To sow\textquotedblright{}.) Say it once more.
\end{tcolorbox}
